% ---------------------------------------------------------------------------
% ADDITION TO THE PERISHABLE-INVENTORY SECTION.
%
% PURPOSE. Extend the perishable results (Table~\ref{tab:perish-results}) with
% three further De Moor et al. (2022) / Farrington et al. (2025) Scenario-A
% instances that the in-crate exact-MDP value iteration also solves exactly
% (each < 2000 states): a longer shelf life (m=3, L=1, FIFO, waste cost 7); the
% first instance with a genuine in-transit pipeline (m=2, L=2, FIFO, waste cost
% 7); and a higher waste cost (m=2, L=1, FIFO, waste cost 10). All three reuse
% EXACTLY the recipe of the two instances already reported (depth-2 oblique-
% linear soft tree, warm-started at the best base-stock gate, CMA-ES under
% common random numbers, model selection on a disjoint validation block; seed
% blocks 64/512/2048, gamma=0.99, horizon 465 with a 100-period warm-up). No
% architecture was re-swept.
%
% INSERTION POINT. Place this block in the perishable section immediately AFTER
% Table~\ref{tab:perish-results} (i.e. right after the \FloatBarrier that
% follows that table) and BEFORE the next section. It introduces one new table,
% \ref{tab:perish-results-more}, and two interpretive sentences.
%
% DATA PROVENANCE (every number sourced from the run artifacts under
% outputs/autoresearch/perishable_inventory_autoresearch/, commit 7b4da10):
%   m=3,L=1,FIFO,cw7  (de_moor2022_m3_exp2_l1_cp7_fifo_d2_oblique_full.json):
%       learned -1425.20, gate -1435.30 (S=8), +0.70%, paired SEM 1.702 ->5.9x,
%       VI analytic -1424.39 (rounded -1424, matches published -1424). 28 params.
%   m=2,L=2,FIFO,cw7  (de_moor2022_m2_exp6_l2_cp7_fifo_d2_oblique_full.json):
%       learned -1462.45, gate -1495.44 (S=9), +2.21%, paired SEM 1.869 ->17.7x,
%       VI analytic -1461.30 (rounded -1461, matches published -1461). 28 params
%       (input includes the one-period in-transit pipeline bucket).
%   m=2,L=1,FIFO,cw10 (de_moor2022_m2_exp4_l1_cp10_fifo_d2_oblique_full.json):
%       learned -1464.06, gate -1485.40 (S=6), +1.44%, paired SEM 1.905 ->11.2x,
%       VI analytic -1463.51 (rounded -1464; published -1463, a one-unit rounding
%       artifact). 21 params.
%   Verdict on all three: beats (gain exceeds the paired SEM). Selection on the
%   disjoint validation block, reported on the held-out 2048-seed eval block.
% ---------------------------------------------------------------------------

\paragraph{Three further exactly-solvable instances.}
To check that the gate beat is not specific to the two anchor instances, we ran the
\emph{identical} recipe --- depth-2 oblique-linear soft tree, warm-started at the best
base-stock gate, CMA-ES under common random numbers, selection on the disjoint validation
block --- on three more De Moor et al.\ Scenario-A settings whose exact-MDP value iteration
the repository also reproduces (each under $2000$ states): a longer shelf life ($m=3$,
$L=1$, FIFO, waste cost $7$), the first setting with a genuine \emph{in-transit pipeline}
($m=2$, $L=2$, FIFO, waste cost $7$), and a higher waste cost ($m=2$, $L=1$, FIFO, waste
cost $10$). No architecture was re-swept; only the longer-state instances enlarge the policy
from $21$ to $28$ parameters, the extra weights coming from the third state component (the
$m=3$ second on-hand bucket and, for $L=2$, the in-transit pipeline bucket). On all three the
learned tree again \emph{beats} the best base-stock gate on the shared Monte-Carlo estimator
by several times its paired standard error (Table~\ref{tab:perish-results-more}), and again
sits essentially on the analytic value-iteration optimum, which we read as context only.

\begin{table}[!htbp]
\centering
\small
\caption{Perishable inventory, three further exactly-solvable instances (companion to
Table~\ref{tab:perish-results}; same recipe and protocol): best learned depth-2 oblique-linear
soft tree vs.\ the best base-stock gate, common-random-number evaluation ($2048$ held-out seeds,
horizon $465$, $\gamma=0.99$). Returns are discounted \emph{negated} costs (higher is better).
$\Delta\%$ is the learned policy's improvement over the gate, ``SEM mult.'' that gain divided by
the paired standard error. The value-iteration optimum is the \emph{analytic} reference of
\citet{deMoor2022perishable} / \citet{farrington2025perishable} and is context only (different
estimator). Instances are described by their parameters; all are FIFO with mean demand $4$,
coefficient of variation $0.5$ and order cap $10$. Source:
\texttt{outputs/autoresearch/perishable\_inventory\_autoresearch/}.}
\label{tab:perish-results-more}
\begin{tabular}{lccccc}
\toprule
Instance ($S$ = best base-stock) & Gate & Learned (ours) & $\Delta\%$ & SEM mult. & VI optimum (context) \\
\midrule
$m=3,\ L=1$, waste $7$ \ ($S=8$)            & $-1435.30$ & $\mathbf{-1425.20}$ & $\mathbf{+0.70\%}$ & $\approx 5.9\times$  & $-1424.39$ \\
$m=2,\ L=2$, waste $7$ \ ($S=9$)            & $-1495.44$ & $\mathbf{-1462.45}$ & $\mathbf{+2.21\%}$ & $\approx 17.7\times$ & $-1461.30$ \\
$m=2,\ L=1$, waste $10$ \ ($S=6$)           & $-1485.40$ & $\mathbf{-1464.06}$ & $\mathbf{+1.44\%}$ & $\approx 11.2\times$ & $-1463.51$ \\
\bottomrule
\end{tabular}
\par\medskip\footnotesize As in Table~\ref{tab:perish-results}, the headline is the like-for-like
gate beat: gate and learned policy are scored on the \emph{same} Monte-Carlo estimator and the gate
is the exact optimum within the base-stock family. The VI-optimum column is the analytic discounted
return under midpoint-binned Gamma demand and is not on the same estimator (it rounds to the published
$-1424$, $-1461$ and $-1463$ respectively), so we report it as context rather than as a separate
improvement claim. Selection is on a disjoint validation block.
\end{table}
\FloatBarrier

\noindent
The pattern of the two anchor instances therefore holds across a longer shelf life, an added lead-time
pipeline, and a higher waste cost: the largest gate beat ($+2.21\%$, $\approx17.7\times$ the paired SEM)
is the $L=2$ pipeline instance, where the age-and-pipeline state gives the learned tree the most room
over a single base-stock level, while the smallest ($+0.70\%$, $\approx5.9\times$) is the $m=3$ case in
which the best base-stock policy is already close to the exact optimum.
